\begin{abstract}
\noindent La bioconversión con larvas de \textit{Hermetia illucens} valoriza residuos orgánicos, pero emite gases de efecto invernadero cuyas fuentes (respiración larvaria y degradación microbiana del lecho) no se han separado experimentalmente en los modelos disponibles. Este trabajo valida un modelo dinámico de emisiones de CO$_2$ construido como extensión del modelo DEB de Eriksen, con tres componentes nuevos: un interruptor de prepupa que cierra la asimilación, un módulo microbiano del lecho y el balance de gas de la cámara estática. El experimento por lotes sellado incluyó dos dietas (alimento estándar y residuos agroindustriales de frutas fermentados), cierres por triplicado (réplicas técnicas), controles de alimento solo y sensores de CO$_2$ y CH$_4$, y la validación explotó la descomposición de fuentes (tratamiento menos control). La componente larvaria reprodujo las tasas netas dentro de la incertidumbre experimental ($R^2 = 0.981$, RMSE de 5.2 ppm/min en la dieta de residuos; en la dieta estándar el colapso de las emisiones se atribuye al agotamiento del sustrato, con una cota inferior incumplida) y el interruptor de prepupa resultó esencial para capturar el colapso de las emisiones al final del ciclo en la dieta de residuos. Usando solo mediciones de gas, el ajuste cuantificó la calidad relativa de la dieta de residuos (70 \% de la estándar), los tiempos de prepupa (12.1 a 16.4 días) y los parámetros microbianos ($k_{\mathrm{ref}} \approx 0.28$ a $0.29$ d$^{-1}$, $Y_{CO_2} \approx 1.13$ a $1.39$); el CH$_4$ se reportó como indicador cualitativo. Como prueba de concepto con un lote por dieta y réplicas técnicas, el modelo y el protocolo de validación son una base para el diseño de dietas y la operación del proceso, y son transferibles a otros insectos de bioconversión.

\vspace{6pt}
\noindent\textbf{Palabras clave:} \textit{Hermetia illucens}; bioconversión; emisiones de CO$_2$; modelo DEB; cámara estática; validación por fuentes.
\end{abstract}
